\reviewsection{Communication, Behavior, and the Reliability of Adult Response}

\citet{peckhamhardin2015behavior} add another important dimension: actions that adults label challenging may be connected to a message that does not yet have a reliable pathway. Functional communication support begins by investigating the purpose, teaching an accessible alternative, and ensuring that partners honor the message. A break message that never produces a break teaches the learner that the new route is less reliable than the action adults are trying to replace.

This builds on my submitted paper, \textit{From Compliance to Communication: Prevent--Teach--Reinforce Planning Through Access, Agency, and Richer Evidence}, without recreating its behavior plan \citep{garciabautista2026ptr}. That paper examined how allowing a familiar laptop to coexist with hands-on materials removed an unnecessary either-or condition. Module 3 adds the language and partner question: how could the learner say, ``keep this,'' ``different order,'' ``show me,'' ``not yet,'' or ``I meant something else'' before an adult completes the interpretation?

Through the Puzzle Plan and EQUITAS, communication support becomes evidence justice. When accessible modes are missing, schools can mistake speech, speed, or compliance for competence. A fair record would document the student's message, available modes, environmental conditions, partner response, breakdowns, successful repairs, family knowledge, and response to support. Even a home--school log can either increase or restrict power. A list of warnings, sitting, following directions, and good/fair/poor ratings tells adults about adult expectations. A student-centered record asks what the learner communicated, what worked, what barrier should change, and what the student wants shared.

This is also the practical test developed in my Module 2 journal, \textit{Who Has More Power After Support?} \citep{garciabautista2026module2journal}. After communication support is introduced, can the student author more messages, enter more relationships, challenge an interpretation, and change what happens next? If visible behavior decreases while adult control remains untouched, communication is still being managed rather than respected.
